% 底本 EGAIV.1.pdf，上传PDF第115--120页，分册印刷第116--121页。
\section{导子与微分}
\source{第四卷第一部分，分册印刷第116--121页；上传PDF第115--120页}
本节引入的概念将在第IV章§16中以几何形式重新讨论，并在预概形的研究中起重要作用。这些概念在本章中的重要性，首先在于它们同形式光滑代数的联系，特别是定理（20.4.9）、（20.5.7）和（20.5.12）。这些定理将在第IV章讨论光滑态射的一节中译成几何语言，在应用中经常使用。另一方面，微分概念还将用于在§22中证明重要的正则性判据；这些判据将在第IV章§7对诺特局部环的深入研究中发挥关键作用。

\subsection{导子与代数扩张}
\egas{20.1.1}{命题}
设 $A$ 为环，不必交换；$B,E,C$ 为 $A$-环，$p:E\to C$ 为核 $\mathfrak a$ 满足 $\mathfrak a^2=0$ 的 $A$-同态，$u:B\to C$ 为 $A$-同态。假设存在 $A$-同态 $v_0:B\to E$ 使 $u=p\circ v_0$。则满足 $p\circ v=u$ 的所有 $A$-同态 $v:B\to E$，恰为形如 $v_0+D$ 的映射，其中 $D:B\to\mathfrak a$ 满足：

$(i)$ $D$ 为 $A$-双模同态；

$(ii)$ 对任意 $f,g\in B$，有
\begin{equation}
 D(fg)=f\cdot D(g)+D(f)\cdot g.
 \tag{20.1.1.1}
\end{equation}
这里 $\mathfrak a$ 通过 $v_0$ 成为 $B$-双模：对 $f\in B$、$z\in\mathfrak a$，令 $f\cdot z=v_0(f)z$、$z\cdot f=zv_0(f)$。
\cmt{I9-N01}{本条没有假设 $p$ 满射，故不把 $\mathfrak a$ 无条件视为整个 $C$ 上的双模。已有提升 $v_0$ 足以定义所需的 $B$-作用；换用另一提升时，差属于 $\mathfrak a$，由 $\mathfrak a^2=0$ 知左右作用不变。当 $p$ 是扩张的增广、因而满射时，才可先定义完整的 $C$-双模结构。}

条件 $p(v(f))=p(v_0(f))$ 正是说 $D(f)=v(f)-v_0(f)\in\mathfrak a$。将 $v(fg)=v(f)v(g)$ 展开，由 $\mathfrak a^2=0$ 得（20.1.1.1）；$(i)$ 则来自 $v,v_0$ 均为 $A$-同态。

设 $\rho:A\to B$ 为结构同态。由（20.1.1.1），对 $a\in A$、$f\in B$ 有
\[
 D(\rho(a)f)=\rho(a)D(f)+D(\rho(a))f.
\]
而 $(i)$ 给出 $D(\rho(a)f)=\rho(a)D(f)$。令 $f=1$，得到 $D(\rho(a))=0$。反过来，若一个\emph{加法映射} $D$ 在 $\rho(A)$ 上为零并满足（20.1.1.1），则它满足 $(i)$。
\cmt{I9-T01}{原文最后一句只写“在 $\rho(A)$ 上为零且满足（20.1.1.1）”，此处将加法性明确写出。Leibniz等式本身不代替加法性；在前面的 $D=v-v_0$ 情形，加法性已经具备。由上述条件，$v_0+D$ 保加法、乘法及结构映射，且 $D(1)=0$，因而确为所需的含幺环同态。}

\egas{20.1.2}{定义}
给定一个 $A$-环 $B$ 和一个 $B$-双模 $L$，满足（20.1.1）的 $(i)$、$(ii)$ 的映射 $D:B\to L$ 称为从 $B$ 到 $L$ 的\emph{$A$-导子}。

由（20.1.1.1）可知，$A$-导子的核是 $B$ 的一个子 $A$-环。有时也将 $\mathbb Z$-导子简称为导子，即满足（20.1.1.1）的加法映射；此时 $D(1)=0$。若 $B$ 是素域 $P$ 上的代数，则每个 $\mathbb Z$-导子都是 $P$-导子。在正特征下，这由上述结论立即得到；在特征零下，对每个非零整数 $n$，由 $n n^{-1}=1$ 得 $nD(n^{-1})=0$，再用 $n$ 在 $P$ 中可逆，得 $D(n^{-1})=0$。

两个 $A$-导子 $D,D':B\to L$ 的差仍为 $A$-导子。因此，所有 $A$-导子组成一个加法群，记为 $\Der_A(B,L)$。若 $A$ 交换，$B$ 为交换 $A$-代数，$L$ 为 $B$-模，则对每个 $b\in B$，$bD$ 仍为 $A$-导子；于是 $\Der_A(B,L)$ 自然成为一个 $B$-模。
\cmt{I9-C01}{依作者正式勘误Err$_{\mathrm{IV}}$10：原文第117页末将 $a\in A$、$aD$ 和“$A$-模”分别改为 $b\in B$、$bD$ 和“$B$-模”。原文的较弱 $A$-模说法并非虚假，勘误补足的是自然的完整标量结构。}

\egas{20.1.3}{推论}
给定 $A$-环之间的两个 $A$-同态 $p:E\to C$、$u:B\to C$，其中 $p$ 的核 $\mathfrak a$ 平方为零。满足 $p\circ v=u$ 的 $A$-同态 $v:B\to E$ 的集合，或者为空，或者是群 $\Der_A(B,\mathfrak a)$ 的一个主齐性空间。

这正是命题（20.1.1）的另一种表述。

\egas{20.1.4}{推论}
设 $A$ 为环，$C$ 为 $A$-环，$L$ 为 $C$-双模，$E$ 为以 $L$ 为核的 $C$ 的一个 $A$-扩张，$p:E\to C$ 为增广。将 $D\in\Der_A(C,L)$ 映为
\[
 x\longmapsto x+D(p(x))
\]
的映射，给出加法群 $\Der_A(C,L)$ 到 $E$ 的自 $A$-等价群的同构。

在（20.1.1）中取 $B=E$、$v_0=1_E$。满足 $p\circ v=p$ 的 $A$-同态恰为 $v=1_E+D'$，其中 $D'\in\Der_A(E,L)$。它是扩张的 $A$-等价，当且仅当它在 $L$ 上等于典范嵌入，即 $D'|_L=0$；这也等价于 $D'$ 可分解为 $E\xrightarrow{p}C\xrightarrow{D}L$，其中 $D$ 是 $A$-导子。由此得到结论。
\cmt{I9-N02}{群运算可直接核对：$pD=0$，故 $(1_E+D_1p)\circ(1_E+D_2p)=1_E+(D_1+D_2)p$；逆为 $1_E-Dp$。这里的等价固定商环 $C$ 和核 $L$，不是任意 $A$-代数自同构。}

\egas{20.1.5}{推论}
设 $A$ 为环，$B,C$ 为 $A$-环，$L$ 为 $C$-双模，$u:B\to C$ 为 $A$-同态。将 $D\in\Der_A(B,L)$ 映为
\[
 x\longmapsto (u(x),D(x))
\]
的映射，给出 $\Der_A(B,L)$ 到所有使 $u$ 分解为 $B\to\mathrm D_C(L)\to C$ 的 $A$-同态 $B\to\mathrm D_C(L)$ 的典范双射；其中 $\mathrm D_C(L)$ 见（18.2.3）。

这是（20.1.3）在平凡扩张上的应用。

\egas{20.1.6}{推论}
设 $A$ 为环，$B$ 为 $A$-环，$L$ 为 $B$-双模。把 $D\in\Der_A(B,L)$ 分别对应到
\[
 (x,y)\longmapsto(x,y+D(x)),\qquad x\longmapsto(x,D(x)),
\]
便在下列三个集合之间建立典范双射：$(i)$ $\Der_A(B,L)$；$(ii)$ 扩张 $\mathrm D_B(L)$ 的自 $A$-等价；$(iii)$ 增广 $\mathrm D_B(L)\to B$ 的所有 $A$-同态右逆。

$(i)$ 与 $(ii)$ 的对应保持群结构；由此诱导的 $(ii)$ 与 $(iii)$ 的对应，就是（18.3.8）中已经定义的对应。

\egas{20.1.7}{推论}
设 $A$ 为环，$C$ 为 $A$-环，$L$ 为 $A$-双模，$E$ 为以 $L$ 为核的 $C$ 的一个平凡 $A$-扩张，$p:E\to C$ 为增广。将 $L$ 赋予该扩张诱导的双模结构。满足 $D(x)=x$（$x\in L$）的 $A$-导子 $D:E\to L$，恰为
\[
 x\longmapsto x-w(p(x)),
\]
其中 $w$ 遍历 $p$ 的所有 $A$-同态右逆。

仍在（20.1.1）中取 $B=E$、$v_0=1_E$。若 $v=1_E-D$，则 $D|_L=1_L$ 等价于 $v|_L=0$，即 $v=w\circ p$，其中 $w:C\to E$ 为 $A$-同态。条件 $p\circ v=p$ 又等价于 $p\circ w=1_C$，也就是 $w$ 为 $p$ 的右逆。

\subsection{导子的函子性质}
\egap{20.2.1}
设 $A$ 为环，$B$ 为 $A$-环，$L,L'$ 为 $B$-双模，$w:L\to L'$ 为 $B$-双模同态。映射 $D\mapsto w\circ D$ 给出加法群同态
\begin{equation}
 w_0:\Der_A(B,L)\longrightarrow\Der_A(B,L').
 \tag{20.2.1.1}
\end{equation}
若 $w':L'\to L''$ 为第二个同态，则 $(w'\circ w)_0=w'_0\circ w_0$。当 $A,B$ 交换且 $L,L'$ 为 $B$-模时，（20.2.1.1）是 $B$-模同态。
\cmt{I9-C02}{作者正式勘误Err$_{\mathrm{IV}}$11把此处的“$A$-模同态”改为“$B$-模同态”，并把本条末尾（20.2.1.3）之后的“关于 $u$ 的模的双同态”改为“$B$-模同态”。下文已按修订文字译入。}

其次，设 $B'$ 为 $A$-环，$v:B'\to B$ 为 $A$-同态，通过 $v$ 将 $L$ 看作 $B'$-双模。则 $D\mapsto D\circ v$ 给出加法群同态
\begin{equation}
 v^0:\Der_A(B,L)\longrightarrow\Der_A(B',L).
 \tag{20.2.1.2}
\end{equation}
这是（20.1.1.1）的结果。对第二个 $A$-同态 $v':B''\to B'$，有 $(v\circ v')^0=v'^0\circ v^0$。若 $A,B,B'$ 交换且 $L$ 为 $B$-模，则（20.2.1.2）是 $A$-模同态。

最后，设 $u:A'\to A$ 为环同态，并据此把 $B$ 看作 $A'$-环。每个 $A$-导子也是 $A'$-导子，故有典范的单射
\begin{equation}
 u^0:\Der_A(B,L)\longrightarrow\Der_{A'}(B,L).
 \tag{20.2.1.3}
\end{equation}
对 $u':A''\to A'$，有 $(u\circ u')^0=u'^0\circ u^0$。当 $A,A',B$ 交换且 $L$ 为 $B$-模时，这也是 $B$-模同态。

于是 $(A,B,L)\mapsto\Der_A(B,L)$ 定义了从（18.3.5）的范畴 $\boldsymbol K$ 到交换群范畴 $\mathbf{Ab}$ 的协变函子：将态射三元组 $(u,v,w)$ 对应到 $w_0\circ v^0\circ u^0$。函子性可由以下两个交换图核对：
\[
\begin{tikzcd}[column sep=large]
 \Der_A(B,L)\arrow[r,"v^0"]\arrow[d,"w_0"']&\Der_A(B',L)\arrow[d,"w_0"]\\
 \Der_A(B,L')\arrow[r,"v^0"']&\Der_A(B',L')
\end{tikzcd}
\]
\[
\begin{tikzcd}[column sep=large]
 \Der_A(B,L)\arrow[r,"u^0"]\arrow[d,"w_0"']&\Der_{A'}(B,L)\arrow[d,"w_0"]\\
 \Der_A(B,L')\arrow[r,"u^0"']&\Der_{A'}(B,L').
\end{tikzcd}
\]

\egas{20.2.2}{定理}
设 $u:A\to B$、$v:B\to C$ 为两个环同态，$L$ 为 $C$-双模。有典范的交换群正合列
\begin{equation}
\begin{aligned}
 &0\longrightarrow\Der_B(C,L)\xrightarrow{u^0}\Der_A(C,L)\xrightarrow{v^0}\Der_A(B,L)\\
 &\qquad\xrightarrow{\partial}\Exan_B(C,L)\xrightarrow{u^1}\Exan_A(C,L)
 \xrightarrow{v^1}\Exan_A(B,L).
\end{aligned}
\tag{20.2.2.1}
\end{equation}
这里 $u^0,v^0$ 分别为（20.2.1.3）、（20.2.1.2）中的同态；$u^1$ 为限制基环的同态，$v^1$ 为扩张的逆像同态，见（18.3.6.4）、（18.3.6.2）。连接同态 $\partial$ 的定义如下：对于 $D\in\Der_A(B,L)$，用 $A$-同态
\[
 \alpha:B\longrightarrow\mathrm D_C(L),\qquad
 x\longmapsto(v(x),D(x))
\]
赋予 $\mathrm D_C(L)$ 一个 $B$-扩张结构，则 $\partial(D)$ 是该扩张的类。这一正合列关于 $L$ 具有函子性，其中所用同态分别由（20.2.1.1）及（18.3.6.1）给出。
\cmt{I9-T02}{原文将扩张函子同态引到（18.3.4.1）、（18.3.3.1）及（18.3.1.1），同实际定义的编号不合。这里明确使用（18.3.6.4）、（18.3.6.2）、（18.3.6.1）；这是译者据源靶核定的引文校订。}

因为 $D$ 是 $A$-导子，上述 $\alpha$ 确为 $A$-同态，因而给出 $B$-扩张，$\partial(D)$ 有定义（20.1.5）。连接映射保持加法：将两扩张先沿 $C$ 取纤维积，再沿 $L\oplus L\to L$ 的加法推出，结构映射的第二分量由 $(D_1,D_2)$ 变为 $D_1+D_2$。
\cmt{I9-N03}{这一句补明原文未展开的 $\partial$ 的加法性，采用§18中扩张类加法的定义。以下五处正合性的顺序仍依原文。}

须在五处核对正合性。

$1)$ 在 $\Der_B(C,L)$ 处，单射性即（20.2.1.3）。

$2)$ $\Ker(v^0)$ 由在 $v(B)$ 上为零的 $A$-导子 $C\to L$ 组成，恰为 $B$-导子，故在 $\Der_A(C,L)$ 处正合。

$3)$ $D\in\Ker\partial$ 等价于由 $\alpha$ 定义的 $B$-扩张是 $B$-平凡的，即存在 $B$-同态截面 $C\to\mathrm D_C(L)$。它当然是 $A$-同态，故由（20.1.6）可写为 $z\mapsto(z,D'(z))$，其中 $D'\in\Der_A(C,L)$。成为 $B$-同态的条件是
\[
 D(x)z+v(x)D'(z)=D'(v(x)z)
 =v(x)D'(z)+D'(v(x))z
\]
对所有 $x\in B$、$z\in C$ 成立；令 $z=1$ 得 $D'(v(x))=D(x)$。因此 $\Ker\partial=\Im v^0$。

$4)$ $\Ker u^1$ 是那些在遗忘为 $A$-扩张后平凡的 $B$-扩张类。把这种扩张与 $\mathrm D_C(L)$ 识别后，其 $A$-结构为
\[
 t\longmapsto(v(u(t)),0).
\]
任何在其上的 $B$-扩张结构均由 $x\mapsto(v(x),D(x))$ 给出，其中 $D$ 是 $\mathbb Z$-导子。其限制为上述 $A$-结构，当且仅当 $D(u(t))=0$，即 $D$ 是 $A$-导子。因此该扩张类正好属于 $\Im\partial$。
\cmt{I9-T03}{原文将这里的 $A$-结构写为 $t\mapsto(u(t),0)$，但 $u(t)\in B$，而第一分量属于 $C$。正文补上复合 $v\circ u$；未改变论证内容。}

$5)$ 一个 $A$-扩张 $E$ 的类属于 $\Ker v^1$，当且仅当存在 $A$-同态 $w:B\to E$，使 $v$ 分解为 $B\xrightarrow{w}E\to C$。这样一个 $w$ 正好在 $E$ 上定义 $B$-扩张结构，其类经 $u^1$ 成为所给类。逆向也成立，故在 $\Exan_A(C,L)$ 处正合。

关于 $L$ 的函子性由上述构造及扩张的推出定义得到。

\egas{20.2.3}{推论}
设 $A,B,C$ 为交换环，$u:A\to B$、$v:B\to C$ 为环同态，$L$ 为 $C$-模。有关于 $L$ 函子性的典范 $A$-模正合列
\begin{equation}
\begin{aligned}
 &0\longrightarrow\Der_B(C,L)\xrightarrow{u^0}\Der_A(C,L)\xrightarrow{v^0}\Der_A(B,L)\\
 &\qquad\xrightarrow{\partial}\Exalcom_B(C,L)\xrightarrow{u^1}\Exalcom_A(C,L)
 \xrightarrow{v^1}\Exalcom_A(B,L).
\end{aligned}
\tag{20.2.3.1}
\end{equation}
证明同（20.2.2），只需核实：对每个 $D\in\Der_A(B,L)$，$\partial(D)$ 所表示的 $B$-扩张是交换 $B$-代数。这由 $C$ 交换及 $L$ 是 $C$-模直接得到。

\egas{20.2.4}{推论}
在（20.2.2）（分别在（20.2.3））的假设下，有典范的、关于 $L$ 函子性的正合列
\begin{equation}
\begin{aligned}
0\longrightarrow\Der_B(C,L)&\longrightarrow\Der_A(C,L)
\longrightarrow\Der_A(B,L)\\
&\longrightarrow\Exan_{B/A}(C,L)\longrightarrow0,
\end{aligned}
\tag{20.2.4.1}
\end{equation}
分别为
\begin{equation}
\begin{aligned}
0\longrightarrow\Der_B(C,L)&\longrightarrow\Der_A(C,L)
\longrightarrow\Der_A(B,L)\\
&\longrightarrow\Exalcom_{B/A}(C,L)\longrightarrow0.
\end{aligned}
\tag{20.2.4.2}
\end{equation}
这是 $\Exan_{B/A}(C,L)$、$\Exalcom_{B/A}(C,L)$ 的定义的结果，见（18.3.7）、（18.4.2）。

\egas{20.2.5}{备注}
假设有交换的环同态图
\[
\begin{tikzcd}
 A\arrow[r]&B\arrow[r]&C\\
 A'\arrow[u]\arrow[r]&B'\arrow[u]\arrow[r]&C'\arrow[u].
\end{tikzcd}
\]
则导子与扩张的上述函子性给出正合列之间的交换图。为使每个对象清楚可读，下面将同一个图沿 $\partial$ 分作两段，$\partial$ 的相容性亦包含在第二段中：
\[
\begin{tikzcd}[column sep=1.2em]
0\arrow[r]&\Der_B(C,L)\arrow[r]\arrow[d]&\Der_A(C,L)\arrow[r]\arrow[d]&\Der_A(B,L)\arrow[d]\\
0\arrow[r]&\Der_{B'}(C',L)\arrow[r]&\Der_{A'}(C',L)\arrow[r]&\Der_{A'}(B',L).
\end{tikzcd}
\]
\[
\begin{tikzcd}[column sep=.8em]
\Der_A(B,L)\arrow[r,"\partial"]\arrow[d]&\Exan_B(C,L)\arrow[r]\arrow[d]&\Exan_A(C,L)\arrow[r]\arrow[d]&\Exan_A(B,L)\arrow[d]\\
\Der_{A'}(B',L)\arrow[r,"\partial"']&\Exan_{B'}(C',L)\arrow[r]&\Exan_{A'}(C',L)\arrow[r]&\Exan_{A'}(B',L).
\end{tikzcd}
\]
对正合列（20.2.3.1）、（20.2.4.1）、（20.2.4.2），也有同样的交换图。
